\section{Uniform Walsh decay for finite-rank operators}\label{sec:walsh}

\subsection{Dimension growth}

We now show that a slow decrease of Walsh norms forces a large Walsh
norm at almost twice the original index. The proof uses the scalar
coefficient lemma to find several independent Walsh factors inside a
single pairing. To make that lemma applicable, we first arrange that
our pairings depend only on their matrices of signs. The initial Walsh
test is indexed by vectors, whereas the scalar lemma concerns arbitrary
sign matrices; the frame construction below removes the extra dependence
on the vectors.

An ordered linearly independent tuple in \(V_n\) will be called a
\emph{frame}. We identify covectors with vectors through the dot
pairing. Independence is required separately on the two sides.

\begin{lemma}[Pairs of frames]\label{lem:frames}
The action
\[
(U,V)\longmapsto(HU,H^{-t}V),\qquad H\in\GL(n,\F_2),
\]
is transitive on pairs of frames of fixed lengths with a prescribed
matrix of pairings \((u_i\cdot v_j)\).
Every binary \(r\times s\) matrix is such a pairing matrix if
\(n\ge r+s\).
\end{lemma}
\begin{proof}
Consider two pairs \((U,V)\), \((U',V')\) with the same pairing
matrix. The covector frames define surjections
\(\pi,\pi':V_n\to\F_2^s\).
The isomorphism \(u_i\mapsto u_i'\) between the primal spans respects
evaluation, and hence maps
\(\Span(U)\cap\ker\pi\) onto \(\Span(U')\cap\ker\pi'\).
Extend this restriction to an isomorphism between the two kernels.
Together with the map on the primal spans, it defines an isomorphism
\[
\Span(U)+\ker\pi\longrightarrow\Span(U')+\ker\pi'
\]
satisfying \(\pi'H=\pi\).
Complete it to an invertible map of \(V_n\) by choosing matching
lifts of a basis complementary to the common evaluation image
\(\pi(\Span(U))=\pi'(\Span(U'))\).
Then \(HU=U'\) and \(\pi'H=\pi\), the latter being equivalent to
\(H^{-t}V=V'\).
For feasibility in dimension \(r+s\), given a matrix \(D\), take
\(u_i=(e_i,0)\) and \(v_j=(D_{:,j},e_j)\).
\end{proof}

\begin{lemma}[Growth from a slow decrease]\label{lem:growth}
Fix an integer \(m\ge1\) and constants \(D_g,D_f\ge1\).
Suppose that \(m\) and the parameters \(a,b\)
in \eqref{eq:lift-parameters} satisfy the hypotheses of
Lemma~\ref{lem:coefficient}. There is \(c_1\in(0,1/2]\), depending only
on these parameters, such that the following holds.
If
\[
c_m(G^*)\le D_g,\qquad c_m(F)\le D_f,\qquad 0\ne T:G\to F,
\]
and \(n>4m\) satisfies
\begin{equation}\label{eq:slow-decrease}
w_{n-2m}(T)\le2w_n(T),
\end{equation}
then
\begin{equation}\label{eq:growth}
w_{2(n-2m)}(T)\ge c_1w_n(T).
\end{equation}
The operator \(T\) need not have finite rank.
\end{lemma}

\begin{proof}
Put \(n'=n-2m\). We will turn a Walsh test at index \(n\) into
bounded scalar functions satisfying Lemma~\ref{lem:coefficient}. A
coefficient of binary rank \(k\ge2\) will then give a Walsh test
at index \(kn'\).

\medskip
\textit{Step 1: realize the Walsh signs exactly and equivariantly.}
Choose families \(p_u\in G\), \(q_v\in F^*\) admissible in
\eqref{eq:walsh-bilinear}, changing a sign if necessary, such that
\[
\gamma=\E_{u,v}\chi(u,v)q_v(Tp_u)\ge w_n(T)/2>0.
\]
Amplify \(T\) over \(V_n\times V_n\), and define
\[
x'_u(\alpha,\beta)=(-1)^{u\cdot\beta+\alpha\cdot\beta}p_{u+\alpha},
\qquad
y'_v(\alpha,\beta)=(-1)^{\alpha\cdot v}q_{v+\beta}.
\]
Uniform translation shows that every \(x'_u\) and \(y'_v\) has
norm at most one in the corresponding finite \(L_2\) space.
Writing \(T'\) for the amplified operator, we obtain
\[
y'_v(T'x'_u)=\gamma\chi(u,v).
\]
Indeed, the pairing includes expectation over \(\alpha,\beta\),
and its phase is
\[
u\cdot\beta+\alpha\cdot\beta+\alpha\cdot v
=(u+\alpha)\cdot(v+\beta)+u\cdot v.
\]

Amplify once more over \(\mathscr G=\GL(n,\F_2)\), obtaining
\(\widetilde T:\widetilde G\to\widetilde F\), and set
\[
x_u(R)=x'_{Ru},\qquad y_v(R)=y'_{R^{-t}v}.
\]
For the bilinear form \(B(x,y)=y(\widetilde T x)\), this gives
\begin{equation}\label{eq:exact-pairing}
x_u\in B_{\widetilde G},\quad y_v\in B_{\widetilde F^*},
\qquad B(x_u,y_v)=\gamma\chi(u,v).
\end{equation}
Use the same right translations
\(\Lambda(H)z(R)=z(RH)\) on the domain and target dual.
They are isometries, preserve \(B\) under simultaneous translation,
and satisfy
\[
\Lambda(H)x_u=x_{Hu},\qquad
\Lambda(H)y_v=y_{H^{-t}v},\qquad
\Lambda(H_1)\Lambda(H_2)=\Lambda(H_1H_2).
\]
The orders follow directly from \((RH)^{-t}=R^{-t}H^{-t}\).
By Lemma~\ref{lem:walsh-basic},
\(w_j(\widetilde T)=w_j(T)\) for every \(j\).
Equation \eqref{eq:cotype-amplification} also preserves the bounds
\(D_g,D_f\), so the moment lifts remain available with the same
\(a,b\).

\medskip
\textit{Step 2: lift frames while preserving the symmetry.}
For every primal \(m\)-frame \(U=(u_i)_{i=1}^m\), choose
\(P_{U,\eps}\in B_{\widetilde G}\) with
\[
\E_\eps\eps_iP_{U,\eps}=ax_{u_i}.
\]
We can make these choices equivariant. For one reference frame,
take the lifts from Lemma~\ref{lem:lifts} and average each of them over
its stabilizer under \(\Lambda\). The prescribed moments are fixed
by this stabilizer, and the averaged lifts are stabilizer-invariant.
Transport them to any other frame by \(\Lambda(H)\).
Two transporters differ on the right by a stabilizer element, so
the result is well-defined and equivariant.

The same construction for covector frames, using the action
\(H^{-t}\), gives \(Z_{V,\delta}\in B_{\widetilde F^*}\) with
\[
\E_\delta\delta_jZ_{V,\delta}=by_{v_j}.
\]
The two systems of lifts can be chosen independently.
Their equivariance, invariance of \(B\), and Lemma~\ref{lem:frames}
show that the following pairings depend only on the displayed sign
patterns and lift labels:
\begin{equation}\label{eq:three-functions}
\begin{aligned}
\Phi(S,\eps,\delta)&=B(P_{U,\eps},Z_{V,\delta}),
& S_{ij}&=\chi(u_i,v_j),\\
\phi(s,\eps)&=B(P_{U,\eps},y_v),
& s_i&=\chi(u_i,v),\quad v\ne0,\\
\psi(t,\delta)&=B(x_u,Z_{V,\delta}),
& t_j&=\chi(u,v_j),\quad u\ne0.
\end{aligned}
\end{equation}
All patterns are feasible since \(n>4m\).
The moment identities and \eqref{eq:exact-pairing} give
\begin{equation}\label{eq:three-moments}
\begin{aligned}
\E_\eps\eps_i\Phi(S,\eps,\delta)&=a\psi(S_{i,:},\delta),
&
\E_\delta\delta_j\Phi(S,\eps,\delta)&=b\phi(S_{:,j},\eps),\\
\E_\eps\eps_i\phi(s,\eps)&=a\gamma s_i,
&
\E_\delta\delta_j\psi(t,\delta)&=b\gamma t_j.
\end{aligned}
\end{equation}

\medskip
\textit{Step 3: control the odd parts on a smaller Walsh cube.}
We need a bound that remains useful after division by \(\gamma\).
For every fixed choice of lift labels, we claim that
\begin{equation}\label{eq:odd-bounds}
\|\Odd_S\Phi(\cdot,\eps,\delta)\|_\infty,\quad
\|\Odd_s\phi(\cdot,\eps)\|_\infty,\quad
\|\Odd_t\psi(\cdot,\delta)\|_\infty
\le w_{n'}(T).
\end{equation}
Decompose \(V_n=V_{n'}\oplus\F_2^m\oplus\F_2^m\), and fix a sign
matrix \(S_0=((-1)^{D_{ij}})\).
For \(u,v\in V_{n'}\), take the primal frame
\[
U(u)_i=(u,e_i,0)
\]
and the two covector frames
\[
V^+(v)_j=(v,D_{:,j},e_j),\qquad
V^-(v)_j=(v,D_{:,j}+\mathbf1,e_j),
\]
where \(\mathbf1\in\F_2^m\) is the all-ones vector.
The tags \(e_i,e_j\) ensure independence for every \(u,v\).
Their sign matrices are \(\chi(u,v)S_0\) and
\(-\chi(u,v)S_0\), so
\[
B\left(P_{U(u),\eps},
\frac{Z_{V^+(v),\delta}-Z_{V^-(v),\delta}}2\right)
=\chi(u,v)(\Odd_S\Phi)(S_0,\eps,\delta).
\]
Both vector families lie in their respective unit balls, and each
depends only on its own Walsh index. Testing
\eqref{eq:walsh-bilinear} at \(n'\), and using amplification
invariance, proves the first bound.

For \(\phi\), retain \(U(u)\) and use the nonzero covectors
\((v,h,e_1)\), \((v,h+\mathbf1,e_1)\), where
\(s_i=(-1)^{h_i}\). Half the difference of the corresponding
\(y\)'s gives the same calculation.
For \(\psi\), use the covector frame \((v,0,e_j)\) and half the
difference of the primal vectors at
\((u,e_1,h)\), \((u,e_1,h+\mathbf1)\).
Their pairings are \(\chi(u,v)(-1)^{h_j}\) and its negative.
The \(e_1\) tags keep the singleton vectors nonzero.
This proves \eqref{eq:odd-bounds}.

\medskip
\textit{Step 4: find a Fourier coefficient of binary rank at least two.}
Define
\begin{equation}\label{eq:normalized-functions}
\begin{split}
f(s)&=\gamma^{-1}\Odd_s\E_\eps\phi(\eps s,\eps),\\
g(t)&=\gamma^{-1}\Odd_t\E_\delta\psi(\delta t,\delta),\\
\Psi(S)&=\gamma^{-1}\Odd_S\E_{\eps,\delta}
                  \Phi(\eps S\delta,\eps,\delta).
\end{split}
\end{equation}
These functions are odd, and the slow-decrease hypothesis gives
\[
\max(\|f\|_\infty,\|g\|_\infty,\|\Psi\|_\infty)
\le w_{n-2m}(T)/\gamma\le4.
\]
Changing variables \(s'=\eps s\) in the first coefficient and using
\eqref{eq:three-moments}, we obtain
\[
\widehat f(e_i)
=\gamma^{-1}\E_{s',\eps}s'_i\eps_i\phi(s',\eps)=a.
\]
The odd projection leaves this coefficient unchanged.
Likewise, \(\widehat g(e_j)=b\).

We next check the row and column projections required by
Lemma~\ref{lem:coefficient}. A first row projection is unchanged by
applying global odd projection, since
\[
\E_\sigma\sigma_i h(-\sigma S)
=-\E_\sigma\sigma_i h(\sigma S).
\]
Substitute \(\xi=\eps\sigma\); then \(\xi,\eps,\delta\) are still
independent uniform sign vectors. By \eqref{eq:three-moments},
\begin{align*}
\gamma\E_\sigma\sigma_i\Psi(\sigma S)
&=\E_{\xi,\eps,\delta}\xi_i\eps_i
                   \Phi(\xi S\delta,\eps,\delta)\\
&=a\,\E_{\xi,\delta}\xi_i
                  \psi(\xi_i(\delta S_{i,:}),\delta)\\
&=a\,\E_\delta\frac{\psi(\delta S_{i,:},\delta)
                         -\psi(-\delta S_{i,:},\delta)}2\\
&=a\gamma\,g(S_{i,:}).
\end{align*}
Interchanging rows and columns, and simultaneously
\((a,\eps,\phi)\) with \((b,\delta,\psi)\), gives
\(\E_\tau\tau_j\Psi(S\tau)=b f(S_{:,j})\).
Thus all hypotheses of Lemma~\ref{lem:coefficient} hold, with the fixed
parameters \(m,a,b\).

Choose a binary matrix \(P\) of rank \(k\ge2\) with
\(|\widehat\Psi(P)|\ge c_*\).
Since \(\Psi\) is odd and this coefficient is nonzero, the total
number of ones in \(P\) is odd. Its coefficient is therefore
unchanged when the odd projection is removed from
\eqref{eq:normalized-functions}.
Writing \(\Phi_{\eps,\delta}(S)=\Phi(S,\eps,\delta)\), row and column
sign multiplication gives the exact identity
\[
\gamma\widehat\Psi(P)
=\E_{\eps,\delta}
\left[\left(\prod_{i,j}(\eps_i\delta_j)^{P_{ij}}\right)
\widehat{\Phi_{\eps,\delta}}(P)\right].
\]
For some fixed \(\eps,\delta\), therefore,
\(h=\Phi_{\eps,\delta}\) satisfies
\begin{equation}\label{eq:h-coefficient}
|\widehat h(P)|\ge c_*\gamma.
\end{equation}
No oddness assertion about this individual function \(h\) is needed.

\medskip
\textit{Step 5: extract the independent Walsh factors.}
Keep \(P,\eps,\delta\) fixed. For arbitrary tuples
\((u'_i)_{i=1}^m\), \((v'_j)_{j=1}^m\) in \(V_{n'}\), set
\[
U_i=(u'_i,e_i,0),\qquad
V(D)_j=(v'_j,D_{:,j},e_j),
\]
where \(D\) ranges over binary \(m\times m\) matrices.
The tagged tuples are frames even if the original tuples have
repeated or zero entries.
The Fourier coefficient formula gives
\begin{equation}\label{eq:coefficient-extraction}
B\left(P_{U,\eps},
\E_D(-1)^{\langle P,D\rangle}Z_{V(D),\delta}\right)
=\widehat h(P)(-1)^{\sum_{i,j}P_{ij}u'_i\cdot v'_j},
\end{equation}
where \(\langle P,D\rangle=\sum_{i,j}P_{ij}D_{ij}\) is computed modulo
two. Indeed, translate the uniform matrix \(D\) by
\((u'_i\cdot v'_j)_{ij}\) in the defining Fourier average.
Both vectors in \eqref{eq:coefficient-extraction} have norm at most
one. The first depends only on the \(u'\)-tuple, and the second only
on the \(v'\)-tuple.

Since \(\rank_{\F_2}P=k\), choose binary \(m\times k\) matrices
\(I,J\) with \(I^tPJ=\Id_k\).
For \(\zeta,\eta\in(V_{n'})^k\), substitute
\[
u'_i=\sum_{\ell=1}^k I_{i\ell}\zeta_\ell,\qquad
v'_j=\sum_{\ell=1}^k J_{j\ell}\eta_\ell.
\]
The exponent in \eqref{eq:coefficient-extraction} becomes
\(\sum_{\ell=1}^k\zeta_\ell\cdot\eta_\ell\).
We have thus obtained a Walsh test on \(V_{kn'}\), of amplitude
\(\widehat h(P)\), with separate unit-ball families on its two sides.
Monotonicity and amplification invariance now yield
\[
w_{2n'}(T)\ge w_{kn'}(T)
=w_{kn'}(\widetilde T)\ge|\widehat h(P)|
\ge(c_*/2)w_n(T).
\]
Taking \(c_1=\min(1/2,c_*/2)\) proves the lemma.
\end{proof}

\subsection{Uniform decay}

The growth lemma turns a sufficiently late slow decrease into an
infinite sequence of slow decreases for the same operator. For a
finite-rank operator, this would contradict its Hilbert-space bound.
This is the only point at which finite rank is used.

\begin{theorem}[Uniform finite-rank decay]\label{thm:uniform-decay}
Fix \(m,D_g,D_f\) satisfying the assumptions of Lemma~\ref{lem:growth}.
There is an integer \(n_0\), depending only on these parameters,
such that every finite-rank \(T:G\to F\) with
\(c_m(G^*)\le D_g\) and \(c_m(F)\le D_f\) satisfies
\begin{equation}\label{eq:uniform-drop}
w_n(T)\le\tfrac12 w_{n-2m}(T)\qquad(n\ge n_0).
\end{equation}
In particular there are numbers \(\rho_N\to0\), independent of
\(G,F,T\), such that
\begin{equation}\label{eq:uniform-decay}
w_N(T)\le\rho_N\|T\|.
\end{equation}
Both bounds also hold for every \(T:G\to F\) that is an
operator-norm limit of finite-rank maps, under the same moment
hypotheses on \(G^*\) and \(F\).
\end{theorem}
\begin{proof}
Let \(c_1\) be given by Lemma~\ref{lem:growth}. Choose an integer
\(K>\log_2(1/c_1)\), and put
\[
n_0=4m(K+2)+2.
\]
Fix a nonzero finite-rank \(T\), and suppose that
\eqref{eq:slow-decrease} holds at some \(n\ge n_0\).
The growth lemma gives
\[
w_{2n-4m}(T)\ge c_1w_n(T).
\]
By monotonicity, the same lower bound holds throughout
\([n,2n-4m]\).

Set \(a_n=\lceil3n/2\rceil\). The \(K\) consecutive blocks
\[
[a_n+2m(j-1),\,a_n+2mj]\qquad(1\le j\le K)
\]
fit in this interval, since
\(2n-a_n=\lfloor n/2\rfloor\ge2mK+4m+1\).
Their endpoint ratios telescope:
\[
\prod_{j=1}^K
\frac{w_{a_n+2m(j-1)}(T)}{w_{a_n+2mj}(T)}
=\frac{w_{a_n}(T)}{w_{a_n+2mK}(T)}\le\frac1{c_1}.
\]
Since \(2^K>1/c_1\), one right endpoint \(t\) satisfies
\[
w_{t-2m}(T)\le2w_t(T),\qquad
t\ge3n/2,\qquad w_t(T)\ge c_1w_n(T).
\]
We can repeat the argument at \(t\).
Every application of the growth lemma returns an inequality for
the original \(T\); its auxiliary amplifications are discarded.
We therefore obtain indices \(n_j\ge(3/2)^j n\) with
\(w_{n_j}(T)\ge c_1^j w_n(T)>0\).
The rank-dependent estimate \eqref{eq:rank-dependent-decay}, for
this one fixed operator, implies
\[
0<w_n(T)\le D_Tc_1^{-j}2^{-n_j/2}\longrightarrow0,
\]
a contradiction.
Thus \eqref{eq:slow-decrease} is impossible for \(n\ge n_0\), proving
the strict form of \eqref{eq:uniform-drop} for nonzero finite-rank
\(T\). The zero operator satisfies the asserted weak inequality.

Iteration proves \eqref{eq:uniform-decay}; for example, take
\[
\rho_N=
\begin{cases}
1,&N<n_0,\\
2^{-\lfloor(N-n_0)/(2m)\rfloor},&N\ge n_0.
\end{cases}
\]
Finally, if finite-rank maps \(T_j\) converge in operator norm to
\(T\), then \eqref{eq:walsh-continuity} gives
\(w_s(T_j)\to w_s(T)\) for every \(s\), while
\(\|T_j\|\to\|T\|\). Passing to the limit proves both estimates
for \(T\).
\end{proof}

\begin{remark}
The constant \(D_T\) may depend on the operator rank. It is used only
to rule out an infinite sequence of slowly decreasing indices for
one fixed operator; the resulting \(n_0\) and \(\rho_N\) depend only
on the fixed moment parameters. The binary ranks of Fourier indices
in Lemmas~\ref{lem:coefficient} and~\ref{lem:growth} play a different role: they
count the independent Walsh factors extracted from a coefficient.
\end{remark}
